\section{Introduction}\label{sec:intro}

We examine piecewise affine (PWA) abstractions for solving range verification problems for feedforward neural networks whose nodes (neurons) feature a wide class of nonlinear and ``learnable'' activation functions.  Range verification involves computing bounds on the outputs of the network, given a set of inputs~\cite{Liu+Others/2021/Algorithms,albarghouthi-book}. Range verification serves as a primitive for verifying properties of neural networks for applications such as control of safety-critical autonomous systems~\cite{Irfan+Others/2020/Towards,katz_barrett_dill_julian_kochenderfer_2017} and robustness of networks to input perturbations~\cite{tjeng_xiao_tedrake_2017,Casadio+Others/2022/NEural}. Commonly used neural networks typically employ activation functions such as sigmoid, ReLU, tanh, Leaky ReLU and GeLU~\cite{DUBEY202292}. These functions are fixed as part of the network architecture and typically do not change as the network is trained. Recently proposed network models such as \emph{deep spline networks}~\cite{Pakshal+Others/2020/Learning,Aziznejad+Others/2020/Deep} and Kolmogorov-Arnold Networks~\cite{li2024kolmogorovarnold,Liu+Others/2025/KAN} employ ``free-form'' activation functions modeled using B-splines or radial basis functions which can change during training time in order to better fit the data.
As a result, such networks have highly nonlinear activation functions  that are not fixed \emph{a priori}. At the same time, commonly used approaches for solving  neural network verification problems have relied on the activation functions of the nodes in the network being piecewise affine (PWA) functions~\cite{katz_barrett_dill_julian_kochenderfer_2017,tjeng_xiao_tedrake_2017} or being abstracted by a PWA function with known error bounds. %~\cite{singh_gehr_püschel_vechev_2019}.
However, the strategy for approximating activation functions by PWA functions is often fixed a priori without consideration of the specific neural network instance.


%T
In this paper, we observe that the choice of a PWA abstraction is, in itself, a decision problem: (a) how many pieces do we use to approximate a given unit in the NN? (b) how does our choice of pieces affect the overall error between the function represented by the NN and its PWA approximation? and (c) how do we optimally choose the number of pieces at each unit locally so that the global error estimate can be minimized?  Our framework proposed in this paper provides answers to these questions through dynamic programming. %The key steps of our approach are illustrated in Fig.~\ref{fig:approach-flow}.


\begin{figure*}[t]
    \centering
    \resizebox{0.6\linewidth}{!}{\begin{tikzpicture}[ node distance = 0.75cm,
  phi/.style={rectangle, draw, text centered },
  summ/.style={circle, draw, fill=red!20, text centered},
  input/.style={circle, draw, fill=black, minimum size=0.3pt, inner sep=0.3pt},
  % Define arrow style
  connection/.style={draw, -latex'}
]

\begin{scope}
    %% First layer approximating each unit (univariate function)
    \node at (-6.5,0.75)[draw=black, fill=yellow!20]{ $1$};
    \begin{scope}[xshift=-4.5cm]
    %% BOUNDING BOX
    \draw (-1.2, -0.2) rectangle (1.2,1.2);
    %% AXES
    \draw [<->, thick] (0, -0.2) -- (0,1.2);
    \draw [<->, thick] (-1.2, 0) -- (1.2,0);
    %% SPLINE 
    \draw[blue, ultra thick] plot [smooth, tension=1] coordinates {(-1,0) (-0.8, 0.2) (-0.5, 0.5) (0.0, 1.0) (0.5, 0.7) (0.7, 0.3) (1.0, 0.0) };
    \end{scope}

    \begin{scope}[xshift=-1cm, yshift =0.5cm]
    %% BOUNDING BOX
    \draw (-1.2, -0.2) rectangle (1.2,1.2);
    %% AXES
    \draw [<->, thick] (0, -0.2) -- (0,1.2);
    \draw [<->, thick] (-1.2, 0) -- (1.2,0);
    %% SPLINE 
    \draw[line width=10pt, draw=blue!10, opacity=0.5] (-1, 0) -- (0.1, 1.0) -- (1,0);
    \draw[blue, ultra thick, dashed] plot [smooth, tension=1] coordinates {(-1,0) (-0.8, 0.2) (-0.5, 0.5) (0.0, 1.0) (0.5, 0.7) (0.7, 0.3) (1.0, 0.0) };

    \draw[thick, red] (-1, 0) -- (0.1, 1.0) -- (1,0);

    \node at (0,-.4) { 2 pieces};
    \node at (0,-.7) { error $= 4.0$};

    \end{scope}

    \begin{scope}[xshift=1.6cm, yshift =0.5cm]
    %% BOUNDING BOX
    \draw (-1.2, -0.2) rectangle (1.2,1.2);
    %% AXES
    \draw [<->, thick] (0, -0.2) -- (0,1.2);
    \draw [<->, thick] (-1.2, 0) -- (1.2,0);
    %% SPLINE 
    \draw[line width=5pt, draw=blue!10, opacity=0.5] (-1, 0) -- (-0.2, 0.9) -- (0.2, 1.0) --  (1,0);
    \draw[blue, ultra thick, dashed] plot [smooth, tension=1] coordinates {(-1,0) (-0.8, 0.2) (-0.5, 0.5) (0.0, 1.0) (0.5, 0.7) (0.7, 0.3) (1.0, 0.0) };

    \draw[thick, red] (-1, 0) -- (-0.2, 0.9) -- (0.2, 1.0) --  (1,0);

    \node at (0,-.4) { 3 pieces};
    \node at (0,-.7) { error $=1.2$};

    \end{scope}

    \begin{scope}[xshift=5cm, yshift =0.5cm]
    %% BOUNDING BOX
    \draw (-1.2, -0.2) rectangle (1.2,1.2);
    %% AXES
    \draw [<->, thick] (0, -0.2) -- (0,1.2);
    \draw [<->, thick] (-1.2, 0) -- (1.2,0);
    %% SPLINE 
    \draw[line width=2pt, draw=blue!10, opacity=0.5]  (-1,0) --  (-0.8, 0.2) -- (-0.5, 0.48) --  (-0.2, 0.9) -- (0.0, 1.0) -- (0.1, 1.03) --  (0.3, 0.9) -- (0.5, 0.7) -- (0.7, 0.3) -- (1.0, 0.0);
    \draw[blue, ultra thick, dashed] plot [smooth, tension=1] coordinates {(-1,0) (-0.8, 0.2) (-0.5, 0.5) (0.0, 1.0) (0.5, 0.7) (0.7, 0.3) (1.0, 0.0) };
    \draw[thick, red] (-1,0) --  (-0.8, 0.2) -- (-0.5, 0.48) --  (-0.2, 0.9) -- (0.0, 1.0) -- (0.1, 1.03) --  (0.3, 0.9) -- (0.5, 0.7) -- (0.7, 0.3) -- (1.0, 0.0);

    \node at (0,-.4) { 10 pieces};
    \node at (0,-.7) { error $=0.02$};

    \end{scope}

    \node at (3.35,0.5) {\large $\cdots$};
    % Calligraphic brace
    \draw [decorate, line width=1pt, decoration = {calligraphic brace}] (-2.5,-0.3) -- (-2.5,1.7);
    \draw [decorate,line width=1pt, decoration = {calligraphic brace}] (6.5,1.7) --  (6.5,-0.3);

    \begin{pgfonlayer}{background}
    \draw[draw=black, thin] (-6.5,2) rectangle (7, -0.5);
    \end{pgfonlayer}
\end{scope}

\begin{scope}[yshift=-3.0cm]
    \node at (-6.5,0.75)[draw=black, fill=yellow!20]{ $2$};

    \node[input](n0i)  at (-6, 2) {}; 
    \node[input, below=1cm of n0i](n1i)  {}; 
    \node[input, below=1cm of n1i](n2i)  {}; 
    \node[input, below=1cm of n2i](n3i)  {}; 

    \node[input, right=0.6cm of n0i](n0)   {}; 
    \node[input, right=0.6cm of n1i](n1)  {}; 
    \node[input, right=0.6cm of n2i ](n2)  {}; 
    \node[input, right=0.6cm of n3i](n3)  {}; 

    \path[->] (n0i) edge node[above]{$x_1$} (n0)
    (n1i) edge node[above]{$x_2$} (n1)
    (n2i) edge node[above]{$x_3$} (n2)
    (n3i) edge node[above]{$x_4$} (n3);

    \node[phi, right=0cm of n0, name=phi11]{\scriptsize $\decor\psi{1}{1}$};
    \node[phi, right=0cm of n1, name=phi12]{\scriptsize $\decor\psi{1}{2}$};
    \node[phi, right=0cm of n2, name=phi13]{\scriptsize $\decor\psi{1}{3}$};
    \node[phi, right=0cm of n3, name=phi14]{\scriptsize $\decor\psi{1}{4}$};

    \node[summ, below right=0cm and 0.6cm  of phi11, name=s1]{$\sum$};
    \node[summ, below right=0cm and 0.6cm  of phi13, name=s2]{$\sum$};

    \path[->] (phi11) edge (s1) 
    (phi12) edge (s1)
    (phi13) edge (s2)
    (phi14) edge (s2); 

    \node[phi, right=0.3cm of s1, name=phi21]{\scriptsize $\decor\psi{2}{1}$};
    \node[phi, right=0.3cm of s2, name=phi22]{\scriptsize $\decor\psi{2}{2}$};
    \node[summ, below right=0.5cm and 0.3cm of phi21, name=s3]{$\sum$};
    \node[phi, right=0.3cm of s3, name=phi31]{\scriptsize $\decor\psi{3}{1}$};
    \node[input, name=o1, right=0.3cm of phi31]{};

    \path[->] (s1) edge (phi21)
    (s2) edge (phi22)
    (phi21) edge (s3)
    (phi22) edge (s3)
    (s3) edge (phi31)
    (phi31) edge node[above]{$y$} (o1);
    \draw [decorate, line width=1pt, decoration = {calligraphic brace}] (0.5,2) --  (0.5,-1);
    \node at (3.7,0.5) { \large $ \begin{array}{rcl} 
    \err & = &  \max_{\vx \in \reals^4} \ | f_N(\vx) - \fhat_N(\vx) | \\ 
    & \leq  & 2.5 \decor\err{1}{1} + 3 \decor\err{1}{2} + \decor\err{1}{3} +  \decor\err{1}{4} \\
    && + 1.5 \decor\err{2}{1} + 0.6 \decor{\err}{2}{2} + 0.2 \decor\err{3}{1} \\ 
    \end{array}$};

    \begin{pgfonlayer}{background}
    \draw[draw=black, thin] (-6.5,-1.5) rectangle (7, 2.5);
    \end{pgfonlayer}
\end{scope}

\begin{scope}[yshift=-6.7cm]
    \node at (-6.5,0.75)[draw=black, fill=yellow!20]{ $3$};
    \node at (-4.3, 1.7){\scriptsize
    $ \begin{array}{|c|ccc|}  
    \hline 
    \# \text{piece} & 1 &  \cdots & 20 \\  
    \decor\err{1}{1} & 5.0 & \cdots & 0.02\\ 
    \hline
    \end{array}$
    };

    \node at (-4.3, 0.9){\scriptsize
    $  \begin{array}{|c|ccc|}  
    \hline 
    \# \text{piece} & 1 &  \cdots & 20 \\  
    \decor\err{1}{2} & 3.2 &  \cdots & 0.01\\ 
    \hline
    \end{array}$
    };

    \node at (-4.3, 0.4){$\cdots$}; 

    \node at (-4.3, -0.1){\scriptsize
    $  \begin{array}{|c|ccc|}  
    \hline 
    \# \text{piece} & 1 &  \cdots & 20 \\  
    \decor\err{3}{1} & 1.7 &  \cdots & 0.1\\ 
    \hline
    \end{array}$
    };

    \draw [decorate,line width=1pt, decoration = {calligraphic brace}] (-2.3,1.8) --  (-2.3,-0.3);    
    \node[input](i0) at (-2.2, 0.75){};

    \node[rectangle, inner sep=5pt, draw=black, right=0.55cm of i0](n0){\begin{tabular}{c}
    \large\textsc{Knapsack}
    \end{tabular}};

    \node[above=0.5cm of n0, name=i1]{$\delta$};
    \node[input, below=0.7cm of n0, name=i2]{};
    \path[->] (i0) edge (n0)
    (i1) edge (n0)
    (n0) edge (i2);

    \node[below=0cm of i2, name=tab1]{
    \scriptsize
    $ \begin{array}{|c|ccc|}  
    \hline 
    \text{node} & \decor{\psi}{1}{1} &   \cdots & \decor\psi{3}{1} \\ 
    \text{\# pieces} & 5 & \cdots & 6 \\ 
    \hline 
    \end{array}$
    };

    \node at (2,0.75)[draw=black, fill=yellow!20]{ $4$};

    \node at (4.5, 0.4){ $\begin{array}{l}
    \max/\min \;\; y \\ 
    \mathsf{s.t} \\
    \; \text{Constraints for PWA model} \\ 
    \; \text{ + approximation error.}
    \end{array}$  };

    \begin{pgfonlayer}{background}
    \draw[draw=black, thin] (-6.5,-1.4) rectangle (1.7, 2.2);
    \draw[draw=black, thin] (2,-1.4) rectangle (7, 2.2);
    \end{pgfonlayer}

\end{scope}
\end{tikzpicture}}
    \caption{An illustration of the key steps of our approach. 1. We consider multiple possible choices for each unit/activation function; 2. We upper bound the overall network-wide approximation error as a linear combination of individual errors; 3. We solve a knapsack problem to choose an optimized approximation that minimizes the error bound; and 4. We connect our PWA abstraction to existing NN verification solvers, especially the mixed integer approach.}
    \label{fig:approach-flow}
\end{figure*}

The key steps of our approach are illustrated in Fig.~\ref{fig:approach-flow}. Let $f_N: \reals^n \rightarrow \reals$ be the function defined by a NN. We assume single output for the ease of presentation. We construct a PWA approximation $\fhat_N: \reals^n \rightarrow \reals$ as follows: 

\begin{compactenum}
\item We consider a set of possible choices for approximating each unit activation function $\psi$ in the network by  PWA functions $\psihat_1, \ldots, \psihat_P$ with increasing number of pieces and decreasing error. We adapt a classic DP algorithm first proposed by Bellman and Roth for this purpose~\cite{Bellman+Roth/1969/Curve,Bellman/1961/Approximation}. We construct a ``trade-off'' table for each unit based on the error $\err_j = \max_{z} | \psi(z) - \psihat_j(z)|$, as the number of pieces $j$  varies.
\item \emph{Error Analysis:} We bound the worst case error $\err_{\max} = \max_{ \vx \in X} |f_N - \fhat_N|$ between the NN and its PWA approximation as a function of the errors made at each individual unit. Theorem~\ref{thm:kan-overall-error} of the paper shows that this error bound can be written as a weighted linear combination of the error bounds obtained for each unit. 
\item \emph{Knapsack Formulation:}  We find a surprising connection between the problem of finding an ``optimized'' abstraction and solving a variant of the \emph{knapsack} problem: a combinatorial optimization problem that involves the optimal choice of items that maximize value while constraining total weights to fall below a budget. 
\item \emph{Connecting Back to Known Activation Functions:} Having identified an optimized allocation of the number of pieces and error at each node of the network, we prove that any PWA function can be written as a weighted combination of ReLU units (Theorem~\ref{thm:relu-encoding}). This allows us to connect our PWA abstraction to a host of existing NN verification techniques that specialize to ReLU network. 
\item \emph{MILP Formulation:} We solve the range verification problem, using  Mixed Integer Linear Programming (MILP), while encoding the error bounds for the PWA approximation  in the MILP model. This allows the verification problem to be translated into a MILP, which can be solved through powerful combinatorial optimization solvers such as Gurobi~\cite{gurobi}, CPLEX~\cite{cplex} and MOSEK~\cite{mosek}. In fact, by translating a PWA function into a weighted sum of ReLU units, our approach may in fact integrate well within existing NN verification tools.


%The result is a sound over-approximation that solves Problem~\ref{def:problem-stmt} with the appropriate tightness guarantee given by the input $\delta$. 
\end{compactenum}

We focus our empirical evaluation on a recently proposed class of feedforward networks with learnable activation functions called Kolmogorov Arnold Networks (KANs)~\cite{li2024kolmogorovarnold,Liu+Others/2025/KAN} that have become popular across many applications in scientific machine learning. Our experiments on KANs trained for four tasks of increasing scale demonstrates that the upfront cost of the knapsack-based abstraction is recovered many times over during verification. \textcolor{red}{Our optimized abstraction yields output bounds that are up to 22$\times$ tighter than uniform PWA allocation, with the largest gap under small-perturbations where verification is most valuable. The advantage grows with both the size of the KAN and the width of the input bounds.} However, we  note that the overhead of computing the optimized abstraction can be  significant, especially for smaller networks, when compared to the time taken to solve the overall range verification problem. That said, the abstraction can be solved once and reused for multiple verification instances. This is often the case for applications such as certifying robustness of neural networks against adversarial perturbations wherein the certification is performed over ranges that arise from a large set of test examples~\cite{SinghEtAl2019AbstractDomain}.
Also, the use of the optimized abstraction shows gains in the computation time for the verification problem in many instances even when the abstraction overhead is taken into account. 

%Additionally, we compare our verification approaches over KANs against established approaches for multi-layer perceptron (MLP) networks for networks trained to similar test set accuracy over a variety of datasets. We note that the approach presented here proves useful in establishing stronger bounds in many cases. Finally, we demonstrate the application of our framework to study the worst case sensitivity of prediction model outputs to tiny perturbations in their inputs. %




